\subsection{The general ambient fibre}

The general fibre has a larger torus action than the one used to
construct the family.  We describe its fan explicitly; this also
determines its singular locus.

\begin{proposition}\label{prop:toricfibre}
For $b\ne0$, the fibre $\cP_b$ is the toric fourfold with the following
primitive rays in $\ZZ\oplus\ker R$:
\begin{equation}\label{eq:x9rays}
\begin{aligned}
r_0&=(-1,-3,-1,0),&r_1&=(-1,-3,0,-1),\\
r_2&=(-1,-1,0,0),&r_3&=(-1,0,0,0),\\
r_4&=(0,1,0,0),&r_5&=(1,-1,0,0),\\
r_6&=(1,3,1,1).
\end{aligned}
\end{equation}
Its maximal cones, written as sets of ray indices, are
\begin{equation}\label{eq:x9maxcones}
\begin{gathered}
1456,\quad1346,\quad0456,\quad0346,\quad0145,\\
0134,\quad01256,\quad1236,\quad0236,\quad0123.
\end{gathered}
\end{equation}
This is a Gorenstein Fano fourfold.
\end{proposition}

\begin{proof}
The vertex levels are $R(v_1)=1$, $R(v_0)=0$, and $R(v_i)=-1$
for $i\ge2$.  With splitting vector $w=v_1$, each nonempty
positive coefficient is its tail cone, with vertex at zero.
A face cone with empty positive coefficient has tail either zero
or $\operatorname{pos}(v_0)$.  Both already occur as positive
coefficients, from $\operatorname{pos}(v_1)$ and
$\operatorname{pos}(v_0,v_1)$ respectively.  Thus the positive slice
is the entire tail fan.  Only the two displaced negative slices can
be nontrivial.  By \cite[Cor.~1.9]{Suess08}, the three-dimensional
torus action extends to a four-dimensional torus action.

Here is the fan, including the gluing data.  Move the two nontrivial
support points to $0$ and $\infty$.  For each marked row of
Table~\ref{tab:x9slices}, take the cone generated by
$(1,A)$, $(-1,B)$ and $(0,\tau)$.  For each unmarked row, take the
two cones generated separately by $(1,A),(0,\tau)$ and
$(-1,B),(0,\tau)$.  Removing proper faces gives
\eqref{eq:x9rays}--\eqref{eq:x9maxcones}.
The sections at first coordinate $1$ and $-1$ recover the two
subdivisions in the table.  Their markings also agree: the first
has eight maximal cells, all marked; the second has ten maximal
cells, eight marked.  The generic tail subdivision has eight
maximal cells, all marked.  Hence these are the marked subdivisions
of the given fibre, which determine the variety
\cite[Rem.~1.6 and Thm.~1.8]{Suess08}.

For completeness, the anticanonical Cartier characters on the ten
cones, in their listed order, are
\begin{equation}\label{eq:x9cartier}
\begin{gathered}
(-2,-1,-2,6),\ (1,-1,-2,3),\ (-2,-1,6,-2),\
(1,-1,3,-2),\ (-2,-1,6,6),\\
(1,-1,3,3),\ (0,1,-2,-2),\ (1,0,-2,0),\
(1,0,0,-2),\ (1,0,0,0).
\end{gathered}
\end{equation}
Each pairs to $-1$ with the rays of its cone and strictly greater
than $-1$ with every other ray.  The seven inequalities
$\langle m,r_i\rangle\ge-1$ define a bounded polytope with exactly
these vertices.  Consequently the displayed cones are its normal
fan, and $-K$ is Cartier and ample.  All the data needed for these
finite comparisons are printed above and in Table~\ref{tab:x9slices};
\path{toric_fibre.sage} reproduces the fan and marked subdivisions.
\end{proof}

\begin{proposition}\label{prop:charts}
The singular locus of $\cP_b$, for $b\ne0$, consists of one
torus-fixed point.
\end{proposition}
\begin{proof}
Nine of the maximal cones in \eqref{eq:x9maxcones} are unimodular.
The remaining cone $01256$ has five extremal rays in dimension
four, with relation
\[
r_0+r_1+r_6=2r_2+r_5.
\]
Its six facets have ray sets
\[
012,\quad015,\quad026,\quad126,\quad056,\quad156.
\]
For each facet, the greatest common divisor of the maximal minors
of its three primitive ray vectors is one.  Thus every proper
face is smooth, while the full cone is singular.  The corresponding
fixed point is therefore the only singular point of the fourfold.
In the original cover, Table~\ref{tab:x9slices} has eight
complete-locus charts and four affine-locus charts.  The only
singular full cone occurs in the complete-locus chart $F_8$;
every lower-dimensional cone arising from an affine-locus chart
is smooth as well.
\end{proof}

\subsection{The relative anticanonical system}\label{sec:h0}

By Proposition~\ref{prop:relativefano}, after shrinking $B$ the
family is proper and Gorenstein and
$\cL=\omega_{\cP/B}^{-1}$ is relatively ample and commutes with
base change.  Constancy and generation of its sections follow
from the special fibre.

\begin{proposition}\label{prop:h0}\label{prop:gg}
After shrinking $B$, the sheaf $E=\pi_*\cL$ is locally free of
rank $162$ and commutes with base change.  The evaluation map
$\pi^*E\to\cL$ is surjective.  In particular
$h^0(\cP_b,\cL_b)=162$ and $\cL_b$ is globally generated for every
$b$ near zero.
\end{proposition}
\begin{proof}
The line bundle $\cL_0$ is ample Cartier on the complete toric
variety $\PP_{\Delta_9}$.  Toric vanishing gives
$H^i(\cP_0,\cL_0)=0$ for $i>0$
\cite[Thm.~9.2.3]{CoxLittleSchenck}, and
$h^0(\cP_0,\cL_0)=\#(\Delta_9^\circ\cap M)=162$.
Cohomology and base change therefore makes $E$ a vector bundle of
rank $162$ commuting with base change
\cite[Tag~07VK]{Stacks}.  Its evaluation map is surjective on
the special fibre, since $\cL_0$ is globally generated.  The
cokernel has closed support whose proper image in $B$ misses
zero.  Removing that image proves the assertion.  After a further
shrinking, trivializing $E$ gives a common $162$-dimensional
space of sections surjecting onto every fibre.
\end{proof}

The general fan provides a direct check of the section count.
Writing a character as $(s,a,b,c)$, its anticanonical inequalities
are
\[
a\ge-1,\qquad a-1\le s\le\min(1,1-a),\qquad
b,c\le1-s-3a,\qquad b+c\ge-1-s-3a.
\]
Only $a=-1,0,1$ occur at lattice points.  At fixed $a,s$ there are
$\binom{5-s-3a}{2}$ choices of $(b,c)$.  Summing gives
\[
(45+36+28+21)+(15+10+6)+1=130+31+1=162.
\]

\subsection{The general anticanonical member}
